\cotefile{34}
\scanpage{34}
\mspage{32}
% archive sheet number: 33
% image line 1
\deleted{Ce \unsure{cas|cor.}, \unsure{est évident} grâce à \unsure{8.6|6.6}.}

7. \uline{Le théorème de monodromie :} \deleted{\unsure{démonstration arithmétique}} \inserted{\uline{Énoncé conjectural}.}

\deleted{L'énoncé est \unsure{le suivant}.}

\uline{Conjecture} 7.1. \inserted{(Serre-Tate)} \unsure{Sous|Soit|Soient}
\deleted{$X$ un schéma de type fini sur $K$, $i$ un entier, $\ell$ un nb premier $\ne p$, \unsure{et $i$ un ent.} \unsure{On suppose} que $H^i(X_{\overline K},\unsure{\mathbb Z/\ell\mathbb Z|Z/\ell Z})$ est fini. Cf. 0.3)}
les conditions de l'énoncé 0.3, si on suppose $\ell\ne p$, la représentation induite sur le groupe d'inertie
\[
 I\longrightarrow \operatorname{Aut}_{\mathbb Z_\ell}\bigl(H^i(X_{\overline K},\mathbb Z_\ell)\bigr)
\]
est quasi-unipotente, i.e. $\exists$ s-groupe d'indice fini $U$ dans $I$ tel que l'image de $U$ est formée d'\unsure{opérations|opérateurs} unipotentes dans $H^i(X_{\overline K}\,\mathbb Z_\ell)$.

7.2. Nous allons voir par la suite deux démonstrations ``relatives'' de cette conjecture, utilisant l'une et l'autre la résolution des singularités, \deleted{\unsure{et les th de pureté}} de façon d'ailleurs assez différente. De ces démonstrations, il résultera que la conjecture est vraie dans chacun des cas suivants :

a) $p=1$, i.e. $k$ (ou encore $S$) de car. nulle.

\margin{\deleted{\unsure{$X$ est propre}}}
b) $i=1$,

c) $S$ est le hensélisé \inserted{(ou le hensélisé strict)} d'un anneau local d'un point d'un schéma de type fini sur un corps, ou sur $\operatorname{Spec}.\mathbb Z$, et $X_K$ est propre sur $K$ \uline{ou} $K$ de car. nulle.

La première démonstration, qui est de nature arithmétique, sera donnée dans le
